\begin{abstract}
Trustworthy AI regulations express normative obligations in natural language, making structured
compliance reasoning across standards challenging. Existing approaches tend to process documents
in isolation, accumulate duplicate facts, and rarely track obligation strength — limiting detection
of conflicts where one standard mandates what another merely recommends.

We introduce two mechanisms absent from prior work: \emph{graph-conditioned novelty
classification} that conditions each step on the current KG, categorizing candidates as
\textsc{partially\_new} or \textsc{new}; and \emph{per-triple deontic modality}, which makes
obligation-strength conflicts detectable. Both are embedded in a seven-stage human-in-the-loop
pipeline with a comparison engine, pre-merge preview, and post-upsert verification.

On MINE-1, the system achieves \textbf{79.4\%} knowledge retention versus 62.7\% for KGGen
(+16.7~points, robust across three LLM judges). Applied to EU AI Act Article~10, ISO/IEC~42001,
and ISO/IEC~5259-3, it detects 41 inter-standard conflicts, 47 normative gaps (26\% of obligated
concepts undefined), and finds 58\% of concepts are addressed by only one standard.
\end{abstract}
